\section{A dynamic finite-experiment consequence}\label{sec:v8-gaussian}
The decision-complete slice of the same spectrum has a different
geometry. We first make its temporal register convention explicit, and
then connect it to the physical count experiment.

Fix $n,r\ge1$, a compact $K\subset\R^r$ with nonempty interior, and a
positive definite covariance $\Gamma$. At successive times observe
independent
\begin{equation}
 Y_t\sim N(\theta_t,\Gamma),\qquad
                 \theta=(\theta_1,\ldots,\theta_n)\in K^n.
                                                        \label{eq:v8-seqgaussian}
\end{equation}
Each $Y_t$ is discarded after the update. No earlier output labels are
retained outside the register. At the terminal time only its final label
is supplied to a parameter-independent reconstruction kernel. This is
not an $S$-symbol channel whose entire transcript is available.

\begin{theorem}[Sequential accumulation under a total register budget]
\label{thm:v8-seqgaussian}
For the cardinality signature with persistent $S$ states and at most
$S$ terminal outcomes, the decision-complete terminal spectrum of
\eqref{eq:v8-seqgaussian} satisfies
\begin{equation}
 c_n S^{-2/(nr)}\le
 \Theta_{G^{(n)}}(S;\mathcal D_{\rm all})
 \le C_n S^{-2/(nr)}.                                 \label{eq:v8-seqgaussianrate}
\end{equation}
The infimum is over causal encoders obeying the total state budget at
\emph{every} time. Lower bounds include arbitrary randomized,
time-dependent encoders. A deterministic causal encoder attains the
upper order. Constants are for fixed $n,r,K,\Gamma$, not uniform in a
growing horizon. The readout/reconstruction randomness is external and
fresh and private, while the physical alphabet and the persistent
register remain finite. No public tape or other side information is
supplied alongside the final label in this theorem.
\end{theorem}
\begin{proof}
The full observation vector is a Gaussian location experiment of
dimension $nr$, with compact mean set $K^n$ and covariance
$I_n\otimes\Gamma$. Every final register gives an experiment with at
most $S$ outcomes, regardless of its causal implementation. The lower
bound of Theorem~\ref{thm:v7-gaussian-budget}, proved below, applies to
all such experiments. Because a register is a garbling of the full
observation, its forward deficiency is zero. Theorem~\ref{thm:v8-duality}
therefore identifies its reverse deficiency with its worst normalized
Bayes regret, proving the lower bound in \eqref{eq:v8-seqgaussianrate}.

For the upper bound put $J=\lfloor S^{1/(nr)}\rfloor$. At time $t$
compactify each coordinate of $Y_t$ by
$\tau(x)=1/2+\arctan(x)/\pi$, quantize into $J$ cells per coordinate,
and append the resulting $r$-tuple of cell indices to the current
register word. At that time there are $J^{rt}\le J^{rn}\le S$
possible words. The known absolute time determines their length; it
carries no data-dependent information. Thus a single $S$-value
register, not $n$ separately uncharged registers, implements the rule.
At the final time apply the positive tensor-hat reconstruction of
Lemma~\ref{lem:v7-hat} in dimension $nr$. For $J\ge2$ its error is
$O(J^{-2})=O(S^{-2/(nr)})$. Enlarge the constant for the remaining
finite budget range. This proves the upper bound using a genuine causal
encoder. Its real-valued input computation has no finite-time claim
under this signature.
\end{proof}

If instead the sequence of emitted cell labels were supplied to the
terminal decoder, its alphabet would be their product. In that different
model a small current register could leave a large data-bearing
transcript. The lower bound above is about the total retained experiment,
not the size of the most recent communication symbol. This is one reason
zero-delay channel coding and total-register coding require separate
resource signatures \cite{LinderYuksel2014,WoodLinderYuksel2017}.

\subsection{A theorem-level downstream arrow}
We now give the adapter in the coordinates of the retained known-mark
exposure protocol, rather than identifying all of the algebraic paper
with this statistical subexperiment.
\begin{corollary}[Known-mark exposure adapter]\label{cor:v8-a2adapter}
Fix positive exposure centre $\lambda^0$, known marks, a derivative
matrix $D$ of full row rank, and $\Lambda=\diag(\lambda^0)$. Put
$\Sigma=D\Lambda D^{\mathsf T}$ and
$B_*=\Lambda D^{\mathsf T}\Sigma^{-1}$. Let $L$ be a fixed
normal-compression map of rank $r\ge1$, let $F$ be an orthonormal
frame for its range, and set
\[
 C=F^{\mathsf T}LD,\qquad A=F^{\mathsf T}L,\qquad
 \Gamma=C\Lambda C^{\mathsf T}.
\]
On a fixed compact index set $H$, observe independent counts of
means $N\lambda_N(h)$, where
$\lambda_N(h)=\lambda^0+N^{-1/2}B_*h$.
If $AH$ contains an $r$-dimensional ball, the deterministic unjittered
statistic $T_{N,J}$ of Section~\ref{sec:v7-a2}, applied to the score
$C(N_\cdot-N\lambda^0)/\sqrt N$, has at most $J^r\le S$ outputs. For
$\mathcal G=\{N(Ah,\Gamma):h\in H\}$, it satisfies
\begin{equation}
 c S^{-2/r}\le\Delta(\mathcal E_{N,J},\mathcal G)
             \le C_K S^{-2/r},\qquad
 J=\lfloor S^{1/r}\rfloor,\quad
 2^r\le S\le K N^{r/6}.                 \label{eq:v8-a2scope}
\end{equation}
Only the physical output cardinality and the stated joint sample/register
window are bounded here. Reconstruction randomization is external,
parameter independent and private; no finite-bit complexity of the raw
real score computation is asserted.
\end{corollary}
\begin{proof}
$DB_*=I$, so $CB_*=A$. Since $F^{\mathsf T}L$ has row rank $r$
and $D$ is surjective, $C$ has row rank $r$ and $\Gamma>0$.
The compact local alternatives have uniformly positive exposures for
all sufficiently large $N$. Thus every hypothesis of
Theorem~\ref{thm:v7-a2} is verified with these actual score weights,
mean set and covariance. That theorem supplies both inequalities in
\eqref{eq:v8-a2scope}; its lower bound allows every experiment with at
most $S$ outcomes, not only this statistic. Theorem~\ref{thm:v8-duality}
gives the all-decision interpretation of the limiting obstruction, and
Theorem~\ref{thm:v8-spectrum} transports each bounded task through the
comparison kernels. No algebraic primary-decomposition theorem is a
premise of this argument.
\end{proof}
The left inequality concerns every experiment with at most $S$
terminal outputs approximating this specified Gaussian family, and
therefore also the physical statistic. The upper inequality concerns
that explicit statistic and its positive reconstruction kernel.
In particular, achieving Gaussian-family distance at most
$C N^{-1/5}$ forces $S\ge cN^{r/10}$, and budgets of that order suffice
for this statistic. These are label-budget statements in the displayed
joint window; they are not a sample-optimality theorem over all
finite-$N$ physical objectives.

Theorem~\ref{thm:v8-duality} supplies a decision interpretation of the
limiting experiment's finite-output obstruction, while
Theorem~\ref{thm:v8-spectrum} transports each bounded task through the
Poisson comparison kernels. These are genuine arrows between specified
experiments. They neither require nor assert a Sinai local limit theorem,
a particle path large-deviation principle, or a physical phase
construction.
